\documentclass[12pt]{article} 
\usepackage{graphicx}
\usepackage[bottom]{footmisc}
\usepackage[utf8]{inputenc} \usepackage[backend=biber, style=authoryear]{biblatex}
\usepackage{imakeidx}
\usepackage[toc,page]{appendix}
\usepackage{listings}
\usepackage{xcolor}
\usepackage{pdflscape}
\usepackage{tabularx}
\usepackage{longtable}

\definecolor{codegreen}{rgb}{0,0.6,0}
\definecolor{codegray}{rgb}{0.5,0.5,0.5}
\definecolor{codepurple}{rgb}{0.58,0,0.82}
\definecolor{backcolour}{rgb}{0.95,0.95,0.92}

\lstdefinestyle{mystyle}{
    backgroundcolor=\color{backcolour},   
    commentstyle=\color{codegreen},
    keywordstyle=\color{magenta},
    numberstyle=\tiny\color{codegray},
    stringstyle=\color{codepurple},
    basicstyle=\ttfamily\footnotesize,
    breakatwhitespace=false,         
    breaklines=true,                 
    captionpos=b,                    
    keepspaces=true,                 
    numbers=left,                    
    numbersep=5pt,                  
    showspaces=false,                
    showstringspaces=false,
    showtabs=false,                  
    tabsize=2
}

\lstset{style=mystyle}
\addbibresource{./bibliography/biblio.bib}
%\linespread{1.6} % delete in final version

\author{Kae Chinchilla 1526925}
\title{Proposal: }
\date{2025-03-05}

%consider using this so it looks better to the eye tbh
%\linespread{1.3}

\begin{document}
\maketitle
\section{Summary}
\section{Introduction} 
In systems programming precise control over data layout is essential
for achieving high performance, interoperability, and correctness. Languages such as C offer minimal 
abstraction over memory, granting programmers the flexibility to write highly efficient code, but also
placing a significant burden on them to manually define binary layouts through
constructs such as bit-fields, unions, and alignment directives. 
While this low-level control is powerful, it is
notoriously error-prone, particularly in contexts where memory correctness is critical, 
such as device drivers, embedded systems, and network protocols.

Traditionally, developers have addressed these challenges by writing manual marshalling and 
unmarshalling code, or by using
unsafe macros to convert between data structures and low-level binary representations. 
These techniques are difficult to maintain, costly in performance, and extremely challenging
to formally verify. As a result, layout-related bugs often lead to serious vulnerabilities,
systems crashes, or subtle performance degradations. A fundamental contributor to this problem
is the lack of clear abstraction between what a data structure represents semantically and how 
it is laid out in memory, complicating both human reasoning and automated analysis.

Real-world examples highlight the seriousness of these issues. In the development
of device drivers, for instance, hardware registers often impose strict alignment
and bit-level layout constraints that must be matched precisely. Misinterpreting
these layouts can lead to misconfiguration of hardware devices, data corruption
or complete system failure. Similarly, in network protocols stacks must
adhere to standardized packet formats, where incorrect struct definitions can result
in vulnerabilities exploitable by malformed input. These risks underline the critical
importance of providing developers with better abstractions for defining and verifying
data layouts. 

Recent research has explored more principled approaches to layout description.
One notable example is \textbf{Dargent}, a data layout
language integrated into the \textbf{Cogent} ecosystem. Dargent enables programmers
to declaratively specify memory layouts for algebraic data types and
supports automatic generation of layout-aware C code with formal correctness
proofs, via integration with the Isabelle/HOL theorem prover. However, 
Dargent is tightly coupled to Cogent, a restricted functional language that,
while powerful for verification purposes, may be difficult to adopt
as it requires programmers to work in a completely different language.

To address this gap, the Cedar project proposes a
data description language designed to operate directly in the C ecosystem.
Cedar enables programmers to define the memory
layout of C structs through a high-level
declarative syntax, while preserving compatibility with C's expressive and 
performance-oriented features. Implemented as a prototype compiler in  Haskell,
Cedar currently supports basic layout declarations. However, its design remains
at an early stage and lacks the robustness and expressiveness required for
practical adoption in systems programming.

This project builds upon the foundational work of Cedar. Its goal is to 
investigate the language and toolchain can be extended to support a 
broader range of C layout idioms, improve tooling robustness, and better
serve the needs of systems programmers. In doing so, it aims to explore the
design principled necessary to develop a sound and practical data description
language for C, balancing flexibility, safety and usability.

\section{Literature Review}
A robust understanding of prior work is essential to motivate and inform
this project, which investigates how to make declarative memory layout
descriptions in C more expressive, practical and safe for systems programming.
This section surveys the current state of the art in memory layout
control, beginning with the challenges posed by low-level programming in C
and the risks of manual layout specification. It then examines existing
techniques and tools used to mitigate these challenges, including runtime
sanitizers and marshalling libraries. Further, the review explores modern systems
programming languages and DSLs, such as Rust, Cap'n Proto, and FlatBuffers, that
aim to improve layout safety and correctness. Finally, the literature review focuses 
on verified layout description approaches, particularly the Dargent system
embedded in Cogent, and Cedar---the subject of this project---which seems to
adapt these ideas to the full C language.

Together, these works outline both the need for and the limitations of current
approaches, and provide a foundation for the investigation carried out in this
research.

\subsection{Systems Programming and Data Layout Challenges}
Systems programming is a field of computer science that focuses
on the design and implementation of low-level software that interacts 
directly with hardware and system resources. This includes operating
systems, device drivers, embedded systems, and other performance-critical
applications. Systems programming is characterized by the need for
high performance, fine-grained control over hardware resources, and efficient use
of limited computational resources. These demands require a deep understanding of
computer architecture, memory management, and low-level programming constructs.


Languages such as C and C++ have traditionally been the primary tools for
systems programming, due to their minimal abstractions over hardware and their
ability to produce highly efficient machine code.
However, this low-level access comes at the
cost of increased complexity and potential for serious errors. In particular, 
the task of managing data layout---determining how structured data is
represented and accessed in memory---poses significant risks. Programmers must
carefully control how structures are packed, how fields are aligned, and how
memory is accessed, often across different platforms with varying architectural
requirements.

Manual control over data layout is critical in many systems applications. For
example, device drivers must interact with hardware registers that demand exact
field alignments and offsets;network stacks must correctly parse protocol
headers with specific bit-level representations; and file systems must
serialize and deserialize data structures according to on-disk formats. Errors
in data layout handling can lead to severe consequences, including memory
corruption, security vulnerabilities, and unpredictable behavior.

Despite the importance of correct data layout, C provides little direct support
for verifying or enforcing these properties. While some compiler-specific 
directives exist (such as \texttt{\#pragma pack} and compiler-specific 
attributes for struct alignment), their use is often fragile, non-portable,
and difficult to reason about formally. Programmers must manually ensure the
correctness of layouts, which increases maintenance costs and the likelihood of
bugs.

These challenges are not merely theoretical. Numerous real-world 
vulnerabilities---such as buffer overflows, misaligned memory accesses, and
serialization bugs---can be traced back to subtle layout issues.
Consequently, there has been significant interest in developing tools,
languages, and frameworks that help programmers specify and reason about
data layouts more safely and declaratively, without sacrificing performance.

\subsection{Manual Approaches and Current Tools}
In response to the challenges of data layout management in C, developers have
traditionally relied on a combination of manual techniques and runtime
debugging tools to mitigate layout-related errors. These approaches
while effective to a degree, remain fundamentally limited in their ability to
ensure correctness, safety, and long-term maintainability.

At the language level, C provides several mechanisms to manually control
memory layout. Programmers can use \texttt{struct} field ordering, bit-fields,
union types, and compiler-specific derivatives such as \texttt{\#pragma pack},
\texttt{\_\_attribute\_\_(packed)}, or alignment hints to influence how data
is laid out in memory. These features are widely used in systems programming,
particularly when interacting with hardware or network protocols that require
binary compatibility with an external specification. 

Despite their widespread use, these techniques are notoriously fragile. 
Their behaviour can vary across compilers and platforms, and small changes in
struct definitions may silently alter the resulting layout due to implicit
padding and alignment adjustments. As a result, layout correctness often
depends on careful inspection and trial-and-error, rather than principled
guarantees.

To aid in the detection of layout-related bugs at runtime, developers commonly employ
dynamic analysis tools such as \textbf{Valgrind} and 
\textbf{AddressSanitizer (ASan)}. Valgrind
performs memory instrumentation to detect invalid memory accesses, buffer
overflows, and uninitialized memory reads. AddressSanitizer, a compiler-based
instrumentation tool, can detect a range of memory errors 
including out-of-bounds accesses and use-after-free bugs. Similarly, 
\textbf{Undefined Behavior Sanitizer (UBSan)} focuses on catching intergers
overflows, type punning violations, and alignment issues. These tools are
invaluable for debugging complex memory errors and are widely used in systems
development and testing pipelines.

While these tools are invaluable for finding certain classes of bugs,
such tools are reactive rather than preventive. They
identify errors only after they occur during runtime, meaning that test coverage
gaps may leave bugs undetected. Moreover, they often impose substantial
performance overhead and may not be suitable for deployment in embedded
or real-time systems. Most critically, they provide no \textit{no compile-time
guarantees} that a memory layout adheres to its intended specification. This
makes them inadequate as a solution to the broader problem of layout correctness,
especially in safety- or security-critical environments.

The most well known language developed for this purpose is Rust, which is a
systems programming language 
that is designed to be safe and efficient. However, when it comes to delicate
data layout, Rust is not always safer, as it may be necessary to use unsafe code
in order to achieve the desired performance. This is where Cider comes in, as it
is a DSL for type definition that is designed to be used in low-level systems
programming, allowing programmers to define types in a safer and more efficient
manner.
\subsection{Dargent and Cogent}
Dargent is a data layout description language and refinement framework integrated
into the Cogent language. Cogent is a purely functional programming language
designed to to generate efficient and verifiable C code for low-level components,
such as file systems and device drivers. It is supported by a certyfying
compiler that generates C code, alongside proofs in Isabelle/HOL.
In this context, Dargent enhances Cogent, by allowing developers to
describe algebraic data types in terms of their memory layout. It aims to
prevent the need for marshalling and un-marshalling data, which is a common
common a common performance bottleneck in low-level programming. Instead, 
Dargent generates C code that directly manipulates the memory along with
setter and getter functions for safe access, all while generating proofs of
correctness.
Dargent's design of allowing layouts to be written independently of types, which
allows the layouts to be evolved independently of the program's logic, all while
ensuring well-formed layouts, ensuring non-overlapping fields, valid offsets,
and bit-aligned representations.
Dargent is an extension of Cogent, and as such it integrates seamlessly into it.
Its compiler generates the getters and setters for each layout, which
manipulate the memory directly. Proof of correctness is also generated during
the compilation process, providing a combination of high-level control
and abstraction with formal guarantees.
\subsection{Cedar}
Cedar builds upon the work done in Dargent, and aims to essnetially be a 
"Dargent for C". It aims to be directly compatible with C, and be more broadly
applicable. It is a layout description language that integrates directly into C,
implemented as a transpiler in Haskell. 
Currently, Cedar allows programmers to memory layout constructs, which are
associated to C structs via syntactic annotations. The transpiler
includes modules for lexical and semantic analysis, reduces the Cedar
data types and layout specifications.
Cedar currently lacks a formal verification backend, however it preserves
the advantages of Dargent: declarative layout specifications, syntactic
clarity, rigorousness. Importantly, as it is designed to work with the
C language, and not a heavily restricted subset, it increases its
practical utility for systems programmers.
The project lays the groundwork for integration with verification systems,
such as CompCert.

CompCert is a formally verified C compiler that provides a certified translation
of C programs into assembly code. It is designed to ensure that the generated
code is semantically equivalent to the original C program, and it provides 
formal proofs of correctness for the translation process. CompCert is based on
the C99 standard and supports a subset of the C language, including most of the
C99 standard library. It is designed to be used in safety-critical systems,
where correctness and reliability are paramount. The decision to use CompCert
as the target for Cedar is motivated by its formal verification capabilities.

\section{Problem Statement and Research Objectives}
How can we extend the work done in Cedar to provide a sound type system for
C that is both expressive and efficient?

\section{Methodology}
This project aims to explore the design of a data description language for C.
In this context we propose the question: \b{How can we design a data description language that is practical
for systems programming, while being sound and expressive for C?}.

\subsection {Design Space Exploration}
This first phase involves a comprehensive literature review to understand the
design space and trade-offs in existing verified systems languages and tools.
This includes:
\begin{itemize}
    \item A review of the tools closest to cider, Dargent and Cogent, to
    understand their desing, verification architecture, layout-type matching
    and formal semantics. 
    \item Reviewing CompCert, as it is the compiler that is being targeted by
    Cider, and to understand how C semantics are formalised. 
    \item A survey of the current state of the art in type systems for systems programming
    in order to identify gaps in the current state of the art.
\end{itemize}
\subsection{Language Design}
This phase of the project involves evaluating the Cider prototype in regards
to its current syntax, semantics and toolings. Drawing from examples in real-world 
use codebases we aim to identify
\begin{itemize}
    \item The limitations of the current implementation
    \item The limitations of the current syntax
    \item The limitations of the current semantics
    \item The limitations of the current tooling
\end{itemize}

\subsection {Implementation and Integration}
The current prototype of Cider will be extended, continuing its codebase written
in Haskell. This will be focusing on modularity and maintainability.
At this stage, the current key development tasks have been identified as:
\begin{itemize}
    \item Refactoring 
\end{itemize}
Where appropriate, small experiments may be conducted to explore optional features.
\subsection {Case studies and Evaluation}
To evaluate the effectiveness of the extended Cedar language, a suite of case
studies will be developed. Given that the project closely follows the work
done in Dargent, similar case studies will be used to evaluate the language:
\begin{itemize}
    \item A structure in power control systems.
    \item A bitfield in a CAN driver
    \item A timer devide drier
    \end{itemize}
These case studies will be used to evaluate the usability and expressiveness
of the language, and to identify any limitations or issues with the implementation
and design. The case studies will be used to evaluate the language in terms
of its usability, expressiveness, and performance. 

\subsection{Iterative Refinement}
Through the project, feedback loops will inform further design choices.
Issues uncovered during case studies or from implementation may lead to 
refinement of syntax, semantics or implementation strategy. Reflections on 
trade-offs between usability, complexity and extensibility will be documented
and discussed in the final report.
\subsection {Optional Extensions}
Aditional test cases may be added to the Cedar test suite, to cover
additional cases that may be found in live codebases. These may include
structures found in the Linux kernel, or other open-source projects.
However, this will depend on the time available, and the
progress made in the previous phases. 

If time allows, we will explore the integration of Cedar with a formal
verification backend. This follows the work done in Dargent, and aims to
provide formal guarantees of correctness for the generated code. Implementation
of this feature would be a significant undertaking, so it will only be
considered if the previous phases are completed ahead of schedule. Additionally,
given the complexity of this task, it may be considered as a separate
project in the future. However, the groundwork for this will be laid
during the project, and the design of the language will be done with
this in mind.
\section {Project Timeline}
\subsection {Potential Impact}
The proposed work has the potential to improve the usability and applicability
of Cedar as a tool for systems programmers who require fine-grained control
over layout without sacrificing readability. By refinning Cedar's language
design and strengthening its implementation, this project aims to contribute
to the broader effort of enabling safer low-level programming practices in C -
a language where layout bugs can have severe consequences for correctness and
performance.



\newpage
\printbibliography

\end{document}
